\section{Results}
\label{sec:results}

\begin{table*}[t]
  \centering \scriptsize
  \setlength{\tabcolsep}{4pt}\renewcommand{\arraystretch}{1.06}
  \caption{Classification accuracy (\%). Mean $\pm$ std over five seeds; $k{=}1$ is the single-checkpoint baseline and shaded bold marks the best $k$ within each dataset block. $^{\dagger}$Tiny ImageNet entries are validation metrics (\cref{tab:datasets}). At $k{=}5$ accuracy beats the baseline in all 30 configurations.}
  \label{tab:results-acc}
  \begin{tabular}{@{}ll cccccc@{}}
    \toprule
    \textbf{Dataset} & $k$ & \textbf{VGG16} & \makecell{\textbf{Dense}\\\textbf{Net121}} & \makecell{\textbf{Efficient}\\\textbf{NetB0}} & \makecell{\textbf{Mobile}\\\textbf{NetV2}} & \makecell{\textbf{ResNet}\\\textbf{101}} & \makecell{\textbf{ViT}\\\textbf{Base}} \\
    \midrule
    \multirow{5}{*}{CIFAR-100} & 1 & \ms{71.57}{0.41} & \ms{75.39}{0.20} & \ms{76.76}{0.51} & \ms{73.23}{0.41} & \ms{75.84}{1.54} & \ms{87.19}{0.65} \\
     & 2 & \ms{72.10}{0.45} & \ms{78.10}{0.20} & \ms{79.85}{0.31} & \ms{74.24}{0.22} & \ms{78.08}{0.28} & \ms{88.68}{0.28} \\
     & 3 & \ms{72.34}{0.40} & \ms{78.73}{0.12} & \ms{80.51}{0.06} & \ms{74.68}{0.24} & \ms{78.60}{0.36} & \ms{89.12}{0.59} \\
     & 4 & \best{72.48}{0.47} & \ms{78.93}{0.18} & \ms{80.92}{0.22} & \ms{74.92}{0.18} & \best{78.82}{0.48} & \ms{89.39}{0.36} \\
     & 5 & \ms{72.46}{0.43} & \best{79.23}{0.30} & \best{81.15}{0.31} & \best{75.06}{0.15} & \ms{78.80}{0.53} & \best{89.64}{0.28} \\
    \midrule
    \multirow{5}{*}{Tiny ImageNet$^{\dagger}$} & 1 & \ms{67.54}{0.85} & \ms{76.82}{0.61} & \ms{80.32}{0.32} & \ms{72.90}{0.54} & \ms{77.50}{0.60} & \ms{89.68}{0.34} \\
     & 2 & \ms{68.04}{0.46} & \ms{77.74}{0.11} & \ms{80.64}{0.25} & \ms{73.39}{0.38} & \ms{80.41}{0.10} & \ms{90.96}{0.15} \\
     & 3 & \ms{68.12}{0.43} & \ms{77.88}{0.13} & \ms{80.72}{0.26} & \ms{73.40}{0.31} & \ms{80.86}{0.13} & \ms{91.31}{0.10} \\
     & 4 & \best{68.35}{0.38} & \best{78.06}{0.15} & \best{80.81}{0.19} & \ms{73.46}{0.23} & \ms{81.23}{0.14} & \ms{91.43}{0.10} \\
     & 5 & \ms{68.34}{0.35} & \ms{78.04}{0.22} & \ms{80.78}{0.20} & \best{73.53}{0.21} & \best{81.34}{0.32} & \best{91.53}{0.12} \\
    \midrule
    \multirow{5}{*}{Burmese Grape} & 1 & \ms{83.76}{2.04} & \ms{84.62}{1.62} & \ms{82.74}{2.09} & \ms{83.89}{1.54} & \ms{85.38}{0.94} & \ms{86.98}{1.92} \\
     & 2 & \ms{85.30}{1.39} & \ms{84.96}{1.34} & \ms{85.56}{1.27} & \ms{85.94}{2.28} & \ms{85.13}{1.49} & \ms{86.58}{2.18} \\
     & 3 & \ms{85.43}{0.84} & \ms{85.00}{1.50} & \ms{85.77}{0.74} & \ms{86.67}{1.53} & \ms{85.17}{1.61} & \best{87.05}{1.44} \\
     & 4 & \ms{85.77}{0.93} & \ms{84.96}{1.40} & \best{86.54}{1.05} & \best{87.19}{1.54} & \best{85.68}{1.41} & \ms{86.79}{1.47} \\
     & 5 & \best{85.94}{0.98} & \best{85.17}{1.50} & \ms{86.45}{1.15} & \ms{87.01}{1.67} & \ms{85.60}{1.28} & \ms{87.01}{1.50} \\
    \midrule
    \multirow{5}{*}{Tomato Leaf} & 1 & \ms{82.08}{2.79} & \ms{84.08}{2.50} & \ms{83.20}{2.46} & \ms{83.04}{3.92} & \ms{81.84}{2.62} & \ms{83.12}{2.93} \\
     & 2 & \ms{83.36}{2.46} & \best{84.56}{2.32} & \ms{84.96}{2.81} & \ms{84.40}{2.34} & \ms{85.36}{1.74} & \ms{83.52}{2.86} \\
     & 3 & \best{84.00}{2.56} & \ms{84.48}{2.66} & \ms{85.20}{2.96} & \ms{85.12}{2.45} & \ms{86.64}{1.94} & \best{83.92}{1.87} \\
     & 4 & \ms{84.00}{2.10} & \ms{84.48}{2.53} & \ms{85.52}{2.55} & \ms{85.92}{3.44} & \best{87.52}{2.57} & \ms{83.68}{2.34} \\
     & 5 & \ms{83.20}{2.68} & \ms{84.40}{2.29} & \best{85.60}{2.30} & \best{86.64}{3.51} & \ms{87.12}{2.56} & \ms{83.68}{2.43} \\
    \midrule
    \multirow{5}{*}{Potato Leaf} & 1 & \ms{86.03}{1.67} & \ms{86.79}{1.88} & \ms{85.94}{0.76} & \ms{87.09}{1.67} & \ms{86.20}{1.02} & \ms{90.30}{1.51} \\
     & 2 & \ms{87.18}{1.08} & \ms{89.10}{1.65} & \ms{88.12}{1.33} & \ms{88.03}{0.80} & \ms{88.42}{1.14} & \ms{90.30}{0.82} \\
     & 3 & \best{88.08}{1.03} & \ms{89.66}{1.34} & \ms{88.46}{0.97} & \ms{88.29}{0.87} & \ms{88.89}{1.42} & \ms{90.43}{0.51} \\
     & 4 & \ms{87.61}{1.38} & \best{90.26}{2.01} & \ms{88.68}{0.74} & \ms{88.50}{1.22} & \ms{88.76}{0.96} & \best{90.90}{0.49} \\
     & 5 & \ms{87.78}{1.26} & \ms{90.00}{2.15} & \best{88.97}{0.69} & \best{88.76}{0.72} & \best{89.10}{0.79} & \ms{90.60}{0.40} \\
    \bottomrule
  \end{tabular}
\end{table*}

\subsection{Classification}

\Cref{tab:results-acc} reports accuracy for the 30 classification
configurations; the corresponding per-configuration macro-F1 table is given in
the supplementary material and is summarized in \cref{tab:aggregate}. At the common setting $k{=}5$ both metrics
improve over the single-checkpoint baseline in \emph{all 30}; at least one $k$
improves over the baseline in all 30 as well, so no configuration is harmed by
adopting the rule.

\paragraph{General-purpose benchmarks.}
On CIFAR-100 the largest CNN gain is EfficientNetB0, $76.76\to81.15\%$
(+4.39 points), while ViT-Base reaches the strongest absolute accuracy,
$87.19\pm0.65\to89.64\pm0.28\%$, with variance more than halved. On Tiny
ImageNet, where entries are validation metrics, ResNet101 gains most,
$77.50\to81.34\%$, and ViT-Base advances $89.68\to91.53\%$ with standard
deviation dropping $0.34\to0.12$.

\paragraph{Agricultural datasets.}
Gains are larger but noisier here, consistent with these being the smallest
datasets. The biggest single improvement anywhere is Tomato Leaf ResNet101,
$81.84\to87.52\%$ accuracy and $84.48\to90.01\%$ macro-F1 at $k{=}4$. On Potato
Leaf, DenseNet121 gains 3.47 points and ViT-Base reaches $90.90\%$ at $k{=}4$.
Burmese Grape is least responsive: MobileNetV2 still gains 3.30 points, but
ViT-Base moves only $86.98\to87.01\%$, the smallest improvement in the study.

\paragraph{Beyond accuracy.}
\Cref{tab:aggregate} aggregates all five recorded classification metrics with
equal weight per configuration; every one increases monotonically in $k$.
Recall gains more than precision (+2.22 vs.\ +1.85 points), so the averaged
model recovers comparatively more of the under-represented classes, which is
why macro-F1 improves slightly more than accuracy. AUC rises only 0.39 points
because the baseline already exceeds 97.7\%, but it rises in 23 of the 24
configurations where it was recorded, so averaging improves the ranking of
predicted scores and not merely the arg-max decision.

\begin{table}[t]
  \centering \small
  \setlength{\tabcolsep}{3pt}\renewcommand{\arraystretch}{1.12}
  \caption{Equally weighted mean over the 30 classification configurations;
    change from the $k{=}1$ baseline in parentheses, in points.
    $^{\ddagger}$AUC averages the 24 configurations for which it was recorded
  (all but Tiny ImageNet).}
  \label{tab:aggregate}
  \resizebox{\linewidth}{!}{%
\begin{tabular}{@{}lccccc@{}}
    \toprule
    $k$ & Accuracy & Macro-F1 & \makecell{Macro\\Prec.} & \makecell{Macro\\Rec.} & AUC$^{\ddagger}$ \\
    \midrule
    1 (base) & 81.73 & 81.65 & 82.58 & 81.58 & 97.77 \\
    2 & 83.10\,{\scriptsize(+1.37)} & 83.19\,{\scriptsize(+1.54)} & 83.94\,{\scriptsize(+1.37)} & 83.06\,{\scriptsize(+1.47)} & 98.05\,{\scriptsize(+0.28)} \\
    3 & 83.48\,{\scriptsize(+1.76)} & 83.53\,{\scriptsize(+1.88)} & 84.17\,{\scriptsize(+1.60)} & 83.44\,{\scriptsize(+1.86)} & 98.12\,{\scriptsize(+0.35)} \\
    4 & 83.72\,{\scriptsize(+1.99)} & 83.76\,{\scriptsize(+2.10)} & 84.38\,{\scriptsize(+1.80)} & 83.70\,{\scriptsize(+2.11)} & 98.15\,{\scriptsize(+0.38)} \\
    5 & \textbf{83.76}\,{\scriptsize(+2.04)} & \textbf{83.86}\,{\scriptsize(+2.21)} & \textbf{84.42}\,{\scriptsize(+1.85)} & \textbf{83.80}\,{\scriptsize(+2.22)} & \textbf{98.16}\,{\scriptsize(+0.39)} \\
    \bottomrule
  \end{tabular}}
\end{table}

\begin{table*}[t]
  \centering \small
  \setlength{\tabcolsep}{6pt}\renewcommand{\arraystretch}{1.08}
  \caption{Detection (PASCAL VOC, box metrics) and instance segmentation (Carparts) results, mean $\pm$ std over five seeds. $k{=}1$ is the single-checkpoint baseline. Shaded bold marks the best $k$ per column within each block. Fitness follows the Ultralytics definition $0.1\,\mathrm{mAP}_{50}+0.9\,\mathrm{mAP}_{50:95}$; for segmentation it sums the box and mask components, which is why it exceeds one.}
  \label{tab:results-det}
  \resizebox{\textwidth}{!}{%
\begin{tabular}{@{}ll cccccc@{}}
    \toprule
    \textbf{Dataset / Model} & $k$ & \textbf{Precision} & \textbf{Recall} & \textbf{mAP$_{50}$} & \textbf{mAP$_{75}$} & \textbf{mAP$_{50:95}$} & \textbf{Fitness} \\
    \midrule
    \multirow{5}{*}{PASCAL VOC / YOLOv8s} & 1 & \ms{0.7779}{0.0059} & \ms{0.7681}{0.0101} & \ms{0.8164}{0.0032} & \ms{0.7067}{0.0028} & \ms{0.6299}{0.0040} & \ms{0.6485}{0.0039} \\
     & 2 & \ms{0.7921}{0.0087} & \ms{0.7872}{0.0049} & \ms{0.8441}{0.0028} & \ms{0.7407}{0.0020} & \ms{0.6611}{0.0027} & \ms{0.6794}{0.0027} \\
     & 3 & \ms{0.7909}{0.0143} & \ms{0.7814}{0.0085} & \ms{0.8385}{0.0153} & \ms{0.7365}{0.0139} & \ms{0.6579}{0.0133} & \ms{0.6760}{0.0135} \\
     & 4 & \ms{0.7920}{0.0078} & \ms{0.7889}{0.0090} & \ms{0.8457}{0.0069} & \ms{0.7434}{0.0042} & \ms{0.6645}{0.0044} & \ms{0.6826}{0.0046} \\
     & 5 & \best{0.7941}{0.0134} & \best{0.7930}{0.0031} & \best{0.8545}{0.0066} & \best{0.7532}{0.0048} & \best{0.6719}{0.0051} & \best{0.6902}{0.0052} \\
    \midrule
    \multirow{5}{*}{PASCAL VOC / YOLO11s} & 1 & \ms{0.7827}{0.0101} & \ms{0.7822}{0.0067} & \ms{0.8317}{0.0063} & \ms{0.7239}{0.0071} & \ms{0.6509}{0.0062} & \ms{0.6689}{0.0062} \\
     & 2 & \ms{0.8114}{0.0124} & \ms{0.8091}{0.0060} & \ms{0.8661}{0.0093} & \ms{0.7647}{0.0099} & \ms{0.6877}{0.0086} & \ms{0.7055}{0.0087} \\
     & 3 & \ms{0.8070}{0.0042} & \ms{0.8056}{0.0063} & \ms{0.8650}{0.0040} & \ms{0.7651}{0.0057} & \ms{0.6884}{0.0050} & \ms{0.7061}{0.0049} \\
     & 4 & \best{0.8155}{0.0075} & \best{0.8119}{0.0030} & \best{0.8766}{0.0020} & \best{0.7778}{0.0038} & \best{0.7001}{0.0036} & \best{0.7177}{0.0034} \\
     & 5 & \ms{0.8127}{0.0139} & \ms{0.8076}{0.0048} & \ms{0.8754}{0.0028} & \ms{0.7740}{0.0047} & \ms{0.6968}{0.0029} & \ms{0.7147}{0.0029} \\
    \midrule
    \multirow{5}{*}{Carparts / YOLOv8s} & 1 & \best{0.6277}{0.0485} & \ms{0.7762}{0.0217} & \ms{0.6941}{0.0157} & \ms{0.6131}{0.0218} & \ms{0.5411}{0.0235} & \ms{1.1361}{0.0415} \\
     & 2 & \ms{0.6183}{0.0309} & \ms{0.8191}{0.0108} & \ms{0.7151}{0.0126} & \ms{0.6467}{0.0183} & \ms{0.5791}{0.0151} & \ms{1.2082}{0.0295} \\
     & 3 & \ms{0.6198}{0.0257} & \ms{0.8194}{0.0208} & \best{0.7166}{0.0161} & \best{0.6523}{0.0212} & \best{0.5826}{0.0165} & \best{1.2152}{0.0342} \\
     & 4 & \ms{0.6103}{0.0313} & \ms{0.8234}{0.0200} & \ms{0.7087}{0.0165} & \ms{0.6443}{0.0232} & \ms{0.5756}{0.0170} & \ms{1.2016}{0.0355} \\
     & 5 & \ms{0.6121}{0.0328} & \best{0.8243}{0.0248} & \ms{0.7087}{0.0173} & \ms{0.6437}{0.0199} & \ms{0.5755}{0.0169} & \ms{1.2025}{0.0353} \\
    \midrule
    \multirow{5}{*}{Carparts / YOLO11s} & 1 & \best{0.6328}{0.0415} & \ms{0.7794}{0.0132} & \ms{0.7066}{0.0267} & \ms{0.6396}{0.0298} & \ms{0.5589}{0.0273} & \ms{1.1706}{0.0538} \\
     & 2 & \ms{0.6211}{0.0143} & \ms{0.8217}{0.0144} & \ms{0.7221}{0.0157} & \ms{0.6648}{0.0131} & \ms{0.5914}{0.0140} & \ms{1.2313}{0.0283} \\
     & 3 & \ms{0.6227}{0.0155} & \ms{0.8269}{0.0290} & \ms{0.7213}{0.0125} & \ms{0.6643}{0.0135} & \ms{0.5925}{0.0112} & \ms{1.2341}{0.0225} \\
     & 4 & \ms{0.6223}{0.0138} & \ms{0.8318}{0.0246} & \ms{0.7260}{0.0113} & \ms{0.6683}{0.0104} & \ms{0.5971}{0.0106} & \ms{1.2452}{0.0207} \\
     & 5 & \ms{0.6209}{0.0205} & \best{0.8386}{0.0262} & \best{0.7263}{0.0089} & \best{0.6708}{0.0072} & \best{0.5984}{0.0084} & \best{1.2471}{0.0159} \\
    \bottomrule
  \end{tabular}}
\end{table*}

\subsection{Detection and Segmentation}

\Cref{tab:results-det} reports the four detection and segmentation
configurations. Fitness---the criterion the pipeline itself optimizes for model
selection---improves at \emph{every} $k$ for all four, gaining 5.8--6.9\%
relative at $k{=}5$: $0.6485\to0.6902$ and $0.6689\to0.7147$ for YOLOv8s and
YOLO11s on VOC, and $1.1361\to1.2025$ and $1.1706\to1.2471$ on Carparts. Every
mAP variant improves in all four configurations, with the strictest threshold
benefiting most on VOC (mAP$_{75}$ +0.050 for YOLO11s), suggesting the averaged
weights improve localization quality and not only detection recall.

\paragraph{A precision--recall trade on Carparts.}
The only two exceptions to uniform improvement are informative. On Carparts,
precision \emph{falls} for both models ($0.6277\to0.6121$ and
$0.6328\to0.6209$) while recall rises far more ($0.7762\to0.8243$ and
$0.7794\to0.8386$), and every mAP metric---which accounts for both---improves.
This mirrors the classification finding that recall gains exceed precision
gains, and is the expected signature of a model that has become less
conservative on rare classes: Carparts has 23 classes over roughly 3.2k
training images. Practitioners for whom precision is the binding constraint
should note this trade.

\subsection{Sensitivity to $k$}
\label{sec:sensitivity}

The best $k$ is configuration-dependent. Over the 30 classification
configurations the highest mean accuracy occurs at $k{=}5$ in 14, $k{=}4$ in
11, $k{=}3$ in four and $k{=}2$ in one; macro-F1 gives 16/10/3/1. Thus
$k\in\{4,5\}$ accounts for 25 of 30. Performance is not uniformly monotonic
per configuration---accuracy is non-decreasing across all of $k{=}1\ldots5$ in
only 12 of 30---because several models saturate at $k{=}3$ or $k{=}4$ and then
dip slightly. Detection and segmentation behave the same way: YOLO11s peaks at
$k{=}4$ on VOC and YOLOv8s at $k{=}3$ on Carparts.

Two things nevertheless hold across the whole study. First, the aggregate is
monotone (\cref{tab:aggregate}), because the per-configuration dips are small
relative to the gains. Second, and more useful in practice, $k{=}5$ improves
the primary metric over the baseline in all 34 configurations even where it is
not the individual optimum. We therefore recommend $k\in\{4,5\}$ as a default
operating range and note that if $k$ is tuned, it must be tuned on the
selection split---never on the evaluation split, as we have done nowhere in
this paper.

\subsection{Run-to-Run Variability}

Averaging usually, but not always, stabilizes the final model. At $k{=}5$ the
across-seed accuracy standard deviation is below the baseline's in 25 of 30
classification configurations and macro-F1 in 23 of 30; averaged over all 30 it
falls from 1.43 to 1.11 points for accuracy ($-22.3\%$) and 1.45 to 1.11 for
macro-F1 ($-23.5\%$), most on the general-purpose benchmarks ($-46\%$
CIFAR-100, $-56\%$ Tiny ImageNet) and least on Tomato Leaf ($-8\%$), the
smallest dataset. Detection and segmentation follow in three of four
configurations, most strikingly YOLO11s on Carparts ($0.0538\to0.0159$,
$-70\%$); YOLOv8s on VOC is the exception, rising $0.0039\to0.0052$. Variance
reduction is a strong tendency, not a guarantee.
